\response Thank you. Fig. 1i-l shows the winter (JJA) sea ice concentration anomalies relative to PI, and the revised Results now describe them. In MH the changes are small. The winter sea ice area grows by 9\% and the ice edge moves north by about 0.3° of latitude. In LIG the ice cover shrinks by 18\%, mainly in the Ross sector (180°–60°W), where the winter ice edge retreats south by 4.7° of latitude. This is the same region as the warm and salty surface anomaly in Fig. 1b,f, and we discuss it further under comment 13. In LGM and MIS3 the ice cover expands strongly. The winter sea ice area grows from 6.0 million km\textsuperscript{2} in PI to 24.3 and 22.0 million km\textsuperscript{2}, and the winter ice edge, averaged around Antarctica, moves north from 65.9°S to 55.0°S and 56.0°S. The largest shift is in the Weddell sector, where the edge moves north by 14–15° of latitude. This expansion is what shields the open ocean from the atmosphere in the glacial states, and it links Fig. 1i-l to the mixed layer changes described in the same paragraph.

The text proposed for the Results reads:

\begin{change}
``The sea ice anomalies (Fig. 1i-l) follow the temperature changes. In MH the winter sea ice area changes little (+9\%). In LIG it decreases by 18\%, mainly in the Ross sector, where the winter ice edge retreats south by about 5° of latitude beneath the warm and salty surface anomaly. In LGM and MIS3 the winter sea ice area grows by a factor of 3.7 to 4.0, from 6.0 million km\textsuperscript{2} in PI to 24.3 and 22.0 million km\textsuperscript{2}, and the ice edge moves north by about 10° of latitude on average and by 14–15° in the Weddell sector.''
\end{change}

\claudenote{revised.tex does not describe Fig. 1i-l yet. The passage above still has to be added to the Results (around line 111). The numbers are 100-year JJA means from the native mesh (fig\_seaice\_obs/seaice\_numbers.txt).}
